\paragraph{F.47 (Proved). Pinned reconstruction over cyclotomic algebraic bases.}
Let $k\subset\overline{\mathbb Q}$ be any subfield contained in
$F\mathbb Q^{\mathrm{ab}}$ for some number field $F$. For all
finite-type $k$-schemes $X,Y$, the natural map
\[
 \operatorname{Hom}_k(X^{\mathrm{awn}},Y)
   \longrightarrow
 \pi_0\operatorname{Hom}^{\mathrm{pin}}_{B G_k}
                 (X_{\mathrm{et}},Y_{\mathrm{et}})
\]
is a bijection, where a relative geometric morphism includes its
commutation isomorphism over $B G_k$. Equivalently, the etale-topos
functor is fully faithful on finite-type schemes after inverting
universal homeomorphisms, with pinned geometric morphisms as the
target morphisms. In characteristic zero the absolute weak
normalization is seminormalization. In particular this settles the
test over $\mathbb Q(\mu_{2^\infty})$ in full, including constants
and noninvertible morphisms, and applies to $\mathbb Q^{\mathrm{ab}}$
and all its finite extensions. We give the derivation from the
verified arithmetic theorems and the reconstruction criterion;
no priority assertion is intended.

\begin{proof}
\subsubsection*{An integrality lemma for cyclotomic algebraic bases}

Let $k\subset\overline{\mathbb Q}$ be contained in
$F\mathbb Q^{\mathrm{ab}}$ for a number field $F$, and let $L/k$ be
finite. Let $b_1,\ldots,b_r\in L$ be distinct and put
$V=\mathbb A^1_L\setminus\{b_1,\ldots,b_r\}$.
Suppose
\[
 \xi\in\varprojlim_n H^1_{\mathrm{et}}(V,\mu_n)
\]
has the following evaluation property: for every $a\in V(L)$,
$a^*\xi$ is the Kummer class of an element of $L^*$.
Then $\xi$ is the Kummer class of an actual unit of $V$.
This is an integrality assertion for a cohomology class; the individual
elements realizing its evaluations need not initially form a rational
function.


Since $\operatorname{Pic}(V)=0$ and
$\mathcal O(V)^*=L^*\times\prod_i(t-b_i)^{\mathbb Z}$, Kummer theory gives
\[
 \xi=\gamma+\sum_{i=1}^r\lambda_i[t-b_i],\qquad
 \gamma\in\widehat{L^*},\quad \lambda_i\in\widehat{\mathbb Z}.
                                                        \tag{F47.R1}
\]
All completions in this argument are $\varprojlim_n A/nA$ (in
multiplicative notation, $\varprojlim_n A/A^n$). We first prove
$\lambda_i\in\mathbb Z$ without assuming that $L$ has a discrete
valuation.

Choose a number field $K_0\subset L$ containing all the $b_i$.
For each $i$, choose a different prime ideal $\mathfrak p_i$ of
$\mathcal O_{K_0}$ such that the $b_j$ are integral there, their pairwise
differences are units there, and the residue field has more than $r$
elements. Only finitely many primes are excluded. The Chinese remainder
theorem supplies $x,y\in K_0$ such that
\[
 v_{\mathfrak p_i}(x-b_i)=1,\qquad
 v_{\mathfrak p_i}(x-b_j)=0\ (j\ne i),\qquad
 v_{\mathfrak p_i}(y-b_j)=0\ (\text{all }j).
                                                        \tag{F47.R2}
\]
For the first condition specify $x$ modulo $\mathfrak p_i^2$; for the
last choose a residue different from each $b_j$. Thus $x,y\in V(L)$.
The elements $r_j=(x-b_j)/(y-b_j)$ are multiplicatively independent
modulo roots of unity: applying the valuations in (F47.R2) to a relation
$\prod_j r_j^{e_j}\in\mu_\infty$ gives $e_i=0$ for every $i$.
This independence persists in every algebraic overfield, since the
relation itself is an equality of the same algebraic numbers.

There is a number field $F_0$ for which $L\subset M=F_0\mathbb
Q^{\mathrm{ab}}$: take $F_0$ to contain $F$ and finitely many algebraic
generators of $L/k$. Theorem~1.1 of Amoroso--Zannier \cite{AmorosoZannier2000} implies
\[
 \bigcap_{n\ge1}(M^*)^n=\mu_\infty.
                                                        \tag{F47.R3}
\]
Indeed $[\mathbb Q^{\mathrm{ab}}(a):\mathbb Q^{\mathrm{ab}}]
\le [M:\mathbb Q^{\mathrm{ab}}]$ for $a\in M$, so their height bound
is uniformly positive on the non-torsion elements of $M$. If
$z_n^n=z\ne0$ for every $n$, then $h(z_n)=h(z)/n$, forcing $z$ to be
a root of unity. Conversely every root of unity is divisible in $M$,
which contains all roots of unity. Tsujimura's Proposition~4.9 \cite{Tsujimura2022}, applied with its parameter $f=1$ and number field
$F_0$, now says that
\[
 A=M^*/\mu_\infty
\]
is a free abelian group. This is the exact specialization of that
theorem used here.

The independent elements $[r_1],\ldots,[r_r]$ of $A$ have finite
support in a basis of the free abelian group. Their coordinate matrix
has rank $r$, so we may choose homomorphisms
$\chi_1,\ldots,\chi_r:A\to\mathbb Z$ such that the integer matrix
$N=(\chi_i([r_j]))_{i,j}$ has nonzero determinant. Choose actual
elements $u_x,u_y\in L^*$ realizing the two evaluations of $\xi$.
Subtracting the two evaluations of (F47.R1), then mapping to
$\widehat A$ and applying the completed characters, gives
\[
 N(\lambda_1,\ldots,\lambda_r)^{\mathrm t}
   =\bigl(\chi_i([u_x/u_y])\bigr)_i\in\mathbb Z^r.
                                                        \tag{F47.R4}
\]
The adjugate matrix shows that
$d\lambda_i\in\mathbb Z$, where $d=\det N\ne0$.
For each prime dividing $d$, membership in $\mathbb Z_p$ forces the
corresponding divisibility of that integer by the prime power in $d$.
Consequently $\lambda_i\in\mathbb Z$. Equivalently,
$\mathbb Q\cap\widehat{\mathbb Z}=\mathbb Z$ inside the finite
adeles. Put $m_i=\lambda_i$.

Finally choose any $a\in V(L)$ and an element $u_a\in L^*$ realizing
$a^*\xi$. The constant class in (F47.R1) is then the Kummer class of
$c=u_a/\prod_i(a-b_i)^{m_i}$. Thus
$\xi=[c\prod_i(t-b_i)^{m_i}]$, as required. When $r=0$ this last
evaluation argument proves the result directly. The actual unit is
unique modulo $\bigcap_n(L^*)^n$, a subgroup of the roots of unity;
the assertion does not identify different pointwise choices before
the translated-function argument is applied.


\subsubsection*{The smooth-curve step in cyclotomic reconstruction}

Fix an algebraic subfield $k\subset\overline{\mathbb Q}$ for which
$k\subset F\mathbb Q^{\mathrm{ab}}$ for some number field $F$.
Every finite extension of $k$ has the same property, after enlarging
$F$. Write $K=\overline{\mathbb Q}$. We use the rational-source
integrality lemma proved above, with all multiplicative completions
taken as $\varprojlim_n A/A^n$.

\emph{Arithmetic input on abelian varieties.}
For each finite extension $L/k$ and each abelian variety $A/L$,
\[
 \bigcap_{n\geq1} nA(L)=0.                                  \tag{F47.1}
\]
Indeed, the finite presentation of $A$ and a finite set of generators
of $L/k$ allow a number field $F_0$ such that
$L\subset F_0\mathbb Q^{\mathrm{ab}}$ and $A$ is obtained by base
change from an abelian variety over $F_0$. Tsujimura's
Theorem~3.1 \cite{Tsujimura2020} says exactly that the Kummer map
on $A(F_0\mathbb Q^{\mathrm{ab}})$ is injective. Its kernel is the
intersection of the multiplication-by-$n$ images, so the injection
$A(L)\hookrightarrow A(F_0\mathbb Q^{\mathrm{ab}})$ proves
(F47.1). The compatible total Tate module of $A(L)$ is also zero:
each coordinate of a compatible torsion system is divisible by every
integer, and (F47.1) kills it. Separately, the height argument in the
rational-source lemma gives
\[
 D(L^*):=\bigcap_n(L^*)^n\subset\mu_\infty.                  \tag{F47.2}
\]
Neither assertion requires that $L$ have a discrete valuation.

\emph{A smooth-curve Kummer specialization lemma.}
Let $C/L$ be a smooth geometrically connected affine curve and let
\[
 \xi\in\varprojlim_n H^1_{\mathrm{et}}(C,\mu_n).
\]
Assume that for every finite $L'/L$ and every $z\in C(L')$ the
specialization $z^*\xi$ is the Kummer class of an actual element of
$(L')^*$. After a finite extension of $L$ there is an actual regular
unit $g$ on $C$ such that $\xi=[g]$. This extension is only an
extension of constants used in the proof; no descent of the given
topos morphism is assumed.

To prove the lemma, let $\overline C$ be the smooth proper completion
and $S=\overline C\setminus C$. Enlarge $L$ so that the points of
$S$ and one point $z_0\in C$ are rational. Purity for a smooth curve
and the Kummer exact sequence give compatible residue maps
\[
 H^1(C,\mu_n)\longrightarrow
       \bigoplus_{s\in S}\mathbb Z/n,
 \qquad
 \operatorname{res}(\xi)=(\lambda_s)_{s\in S},
 \quad\lambda_s\in\widehat{\mathbb Z}.                    \tag{F47.3}
\]
We first prove that every $\lambda_s$ is an integer.

Fix $s\in S$. By Riemann--Roch and interpolation of values and a
first derivative at finitely many points, there is a nonconstant
rational function
\[
 \pi:\overline C\longrightarrow\mathbb P^1_L
\]
which has a simple zero at $s$ and takes nonzero values at every
other point of $S$. For example allow a sufficiently large pole
divisor at a point different from $S$, prescribe the first-order
jet of a uniformizer at $s$, and prescribe the value $1$ at the
other boundary points. The map is finite. Remove from $\mathbb P^1$
the values $0,\infty$, its branch values, and $\pi(S)$, and enlarge
$L$ once more to split this finite set if necessary. This gives
\[
 W=\pi^{-1}(V)\subset C,\qquad W\longrightarrow V
 \quad\hbox{finite etale},
\]
where $V$ is a punctured affine line. Apply the compatible norm
maps to obtain
\[
 \eta=\operatorname{Cor}_{W/V}(\xi|_W)
        \in\varprojlim_n H^1(V,\mu_n).
\]
The specialization property is preserved. If a fibre above
$a\in V(L')$ is the disjoint union of $\operatorname{Spec}L_i$,
etale base change for the norm and the usual Kummer norm identity
give the Kummer class of
\[
        \prod_i N_{L_i/L'}(u_i)\in(L')^*,
\]
where $u_i$ realizes the evaluation of $\xi$ at the corresponding
point of $W$. Thus the rational-source integrality lemma applies to
$\eta$.

The residue at $0$ of a corestriction is
$\sum_{z\mid0}[\kappa(z):L]\operatorname{res}_z(\xi)$.
This is the divisor-of-norm formula, and can also be read directly
from the norm of a local parameter in the Kummer boundary map.
Points of $C$ contribute zero because $\xi$ extends over them;
among the points of $S$ only the rational point $s$ lies over $0$.
Consequently this residue is exactly $\lambda_s$, with coefficient
one. The lemma on $V$ makes it integral. The choice of a simple
zero also makes explicit that no ramification factor has been
silently inserted.

Put $D=\sum_s\lambda_s[s]$, an integral divisor. The next arrow in
the localization sequence sends the residue vector to the Kummer
Chern class of $D$. Therefore
\[
 c_1(D)=0\quad\hbox{in }H^2(\overline C,\mu_n)
       \quad\hbox{for every }n,
\]
and the Kummer injection
$\operatorname{Pic}(\overline C)/n\hookrightarrow
H^2(\overline C,\mu_n)$ implies
$[D]\in\bigcap_n n\operatorname{Pic}(\overline C)$.
Its degree is zero. Any divisor class $E$ with $nE=[D]$ likewise
has degree zero. The rational point $z_0$ identifies degree-zero
line bundles with $J(L)$ for the Jacobian $J$ and removes the
Brauer obstruction to representing a Picard-scheme point by a line
bundle. Equation (F47.1) therefore gives $[D]=0$.
Choose $g_0\in L(C)^*$ with $\operatorname{div}(g_0)=D$;
it is a unit on $C$.

All residues of $\xi-[g_0]$ now vanish. For every $n$ the map
$H^1(\overline C,\mu_n)\to H^1(C,\mu_n)$ is injective: purity
identifies its preceding support group with a negative-degree
cohomology group of $S$. Thus its residue-free lifts are unique,
and hence compatible as $n$ varies. The Kummer sequences give
\[
 0\longrightarrow\widehat{L^*}\longrightarrow
  \varprojlim_n H^1(\overline C,\mu_n)
  \longrightarrow T\operatorname{Pic}(\overline C)
  \longrightarrow0.
\]
Here the left inverse system has surjective transition maps, so
there is no $\varprojlim^1$ obstruction. The total Tate module on
the right is zero by (F47.1). Thus $\xi-[g_0]$ is a constant
completed Kummer class. Evaluate at $z_0$: the hypothesis identifies
it with the Kummer class of an actual constant $c\in L^*$.
Then $g=cg_0$ proves the lemma.

\emph{Algebraicity of the point function for a smooth curve.}
Let
$\phi:C_{\mathrm{et}}\to(\mathbb G_m)_{\mathrm{et}}$
be pinned and relative to $B G_L$, and write
$f:C(K)\to K^*$ for its point function. It is Galois equivariant
and sends $L'$-points to $L'$-points for every finite extension,
as checked in F.46. Its pulled-back Kummer class satisfies the
specialization lemma. After a finite extension choose
$g_0\in\mathcal O(C)^*$ representing it. Equation (F47.2) gives
\[
 f(x)/g_0(x)\in\mu_\infty\qquad(x\in C(K)).                 \tag{F47.4}
\]
If the generic image is a closed point, the inverse-image of its
open complement is empty by continuity and the Jacobson property.
The morphism factors through that closed subtopos. After splitting
its residue field it has a constant point function. We may therefore
assume the generic image is generic.

For each $b\in K$, base change to contain $b$ and pull back the
open $\mathbb G_m\setminus\{b\}$. It is a dense open
$C_b\subset C$ whose omitted geometric points are finite, and
$x\in C_b(K)$ exactly when $f(x)\ne b$. Applying the curve lemma
to the Kummer classes of $t-b$ gives a rational function $g_b$, a
unit on $C_b$, with
\[
 f(x)-b=d_b(x)g_b(x),\qquad d_b(x)\in\mu_\infty,
       \quad x\in C_b(K).                                 \tag{F47.5}
\]
Use the fixed proper field $M=F\mathbb Q^{\mathrm{ab}}\subset K$.
In any one argument, a finite extension $N/M$ contains the finitely
many coefficients needed there; the field $M$ itself stays fixed.
All the coefficients
$d_b(x)$ lie in $M^*$. The extension $K/M$ has infinite degree,
as does $K/N$ for every finite extension $N/M$ in $K$.
For clarity, the latter assertion can also be seen from the height
bound: if $K=N$ were finite over $M$, that bound would be uniformly
positive on its nontorsion elements, whereas $K$ contains all roots
of the rational number $2$.

The finite-exception quasi-rationality argument in F.46 applies to
an arbitrary smooth curve, with the following explicit replacement
for its polynomial-degree step. Descend the curve and the finitely
many rational functions under consideration to a finite extension
$N/M$. The poles of any fixed linear combination of two such
functions lie in one fixed divisor. If the combination is nonzero,
each zero has residue degree over $N$ bounded by the degree of that
divisor. Points of arbitrarily large residue degree exist: choose a
finite map of the proper curve to $\mathbb P^1_N$, choose values of
arbitrarily large degree over $N$, and take points above them.
They avoid any prescribed finite exceptional set. Thus a vanishing
at such a point forces the same rational-function identity as in
F.46. The two shifted outside-$M$ sets cover the cofinite domain,
and have infinite intersection by this same degree bound.

If $g_0$ is nonconstant, this proves that $f$ equals a rational
function $R\in K(C)$ on a cofinite subset. If $g_0$ is constant,
divide $f$ by one of its values and use (F47.4) to put its values
in $\mu_\infty\subset M$. Choose $b\notin M$. If $g_b$ were
nonconstant, (F47.5) would put its cofinite image in the vector
space $M+Mb$, which cannot contain a cofinite subset of $K$.
Hence $g_b$ is constant. Two different values of $f$ would, by
subtracting their equations, force this constant into $M$, and
then force $b\in M$. Thus $f$ is constant in this case.

We remove the finite exceptional set in the nonconstant case.
For every $b$, the rational function $(R-b)/g_b$ takes its values
in $\mu_\infty$ on a cofinite set. A nonconstant rational function
on a smooth curve has cofinite image in $K$, so this quotient is a
constant $c_b\in\mu_\infty$. Fix $x\in C(K)$ and choose
$b\ne f(x)$. Since $g_b$ is a unit at $x$, the identity
$R-b=c_b g_b$ first excludes a pole of $R$ at $x$. For every
$b\notin\{f(x),R(x)\}$, it then gives
\[
 \frac{f(x)-b}{R(x)-b}\in\mu_\infty.
\]
If $f(x)\ne R(x)$ this fractional linear expression takes all
values of $K\setminus\{0,1\}$ as $b$ varies, a contradiction.
Therefore $f=R$ on every geometric point. It follows that $R$ is
a regular unit on $C_K$. Galois equivariance and equality on the
Zariski-dense geometric points imply that $R$ descends to a regular
unit on the original curve over $L$. This proves the full
point-function assertion for every smooth affine curve and every
finite extension of the base field.

It also gives the equality of compatible Kummer classes that will
be used below. Over the finite extension $L'/L$ employed in the
specialization lemma, the rational function $R/g_0$ has values in
$\mu_\infty$ at every geometric point and is therefore a constant
$\zeta\in(L')^*$. Evaluation at a rational point of the curve gives
$\zeta\in D((L')^*)$, so $[R]=[g_0]=\xi$ over $L'$.
If $C(L)$ is nonempty, equality already holds over $L$. Indeed,
for every $n$ the residual class in $H^1(C,\mu_n)$ is geometrically
zero, and hence is a constant class by the low-degree
Hochschild--Serre sequence. Its evaluation at an $L$-point is zero,
because $R$ is the actual point function there. The same sequence,
or the rational-point section, shows that this constant class is
zero. Without an $L$-point, equality after a finite extension is all
that is asserted or needed for the geometric residue tests.

\subsubsection*{From the curve calculation to arbitrary dimension}

We give the higher-dimensional step independently, so that the passage
from point values to a geometric morphism does not rest on a
fundamental-group calculation. Let $k\subset\overline{\mathbb Q}$ be
contained in $F\mathbb Q^{\mathrm{ab}}$ for some number field $F$.
We use the preceding curve calculation in precisely the following
form: after every finite extension $L/k$, a pinned relative geometric
morphism $C_{\mathrm{et}}\to(\mathbb G_{m,L})_{\mathrm{et}}$, with
$C$ a smooth geometrically connected affine curve having an $L$-point,
has the geometric-point values of a unit of $C$. Its compatible Kummer
class is the class of that unit. The conclusion below is conditional
only on that already proved curve calculation and the arithmetic
inputs explicitly recorded here.

It is enough that the curve calculation initially represents its
Kummer class after a further finite extension. Indeed, suppose the
pointwise unit is $f_C\in\mathcal O(C)^*$ and, over a finite
extension $L'/L$, the Kummer class is $[g_0]$. At every algebraic
point the value of $f_C/g_0$ is divisible in the multiplicative group
of a finite extension of $L'$ and hence is a root of unity, by the
height input used in the curve calculation. A nonconstant rational
function on a geometrically integral curve has cofinite image in
$\mathbb G_m(\overline L)$, so this ratio is constant. Evaluation at
an $L$-point shows that its constant value is itself divisible in
$L'^*$; hence its Kummer class is zero. Thus the two classes agree
after $L'$. At each finite Kummer level, their difference on $C$
lies in the image of $H^1(L,\mu_n)$ by the Hochschild--Serre exact
sequence, since it is zero geometrically. Evaluation at the $L$-point
is a retraction onto this image and is zero on the difference.
Consequently the classes already agree over $L$. This proves exactly
the curve input used below, without requiring a chosen function to
descend merely because its Kummer class descends.

\emph{The arithmetic Picard lemma.}
For every finite extension $L/k$ and every abelian variety $A/L$,
\[
                 \bigcap_{n\geq1}nA(L)=0.                 \tag{F47.H1}
\]
Indeed, choose a number field $F_0$ over which $A$ is defined and large
enough that $L\subset M=F_0\mathbb Q^{\mathrm{ab}}$. Theorem~3.1 of
Tsujimura (2020) states that the Kummer map on $A(M)$ is injective
\cite{Tsujimura2020}. Its proof explicitly factors this map as
\[
 A(M)\longrightarrow\varprojlim_n A(M)/nA(M)
       \lhook\joinrel\longrightarrow H^1(G_M,T A),
\]
so its kernel is exactly $\bigcap_n nA(M)$. Thus (F47.H1) follows by
inclusion. Equivalently, the finite nonempty division fibres of a point
in $\bigcap_n nA(M)$ have a compatible inverse-limit point: use the
factorial levels and the finite intersection property. There is no
distinction here between unrelated divisions and coherent divisions.

If $Y/L$ is smooth, projective, and geometrically integral, put
$P=\operatorname{Pic}(Y)$ and let $P^0$ be the subgroup whose geometric
classes are algebraically trivial. Then
\[
                 \bigcap_{n\geq1}nP=0,
                 \qquad T P=\varprojlim_n P[n]=0.         \tag{F47.H2}
\]
Here $P^0$ injects into the group of $L$-points of the abelian variety
$\operatorname{Pic}^0_{Y/L}$, and
$N=P/P^0$ is a subgroup of the finitely generated geometric
Neron--Severi group; for the latter finiteness statement see
\cite[Theorem~7.3 and Lemma~7.2]{FA_ChengEtAl2021}.
These assertions do not require the Picard scheme
to represent every rational Picard class: the possible Brauer
obstruction only changes the image of $P^0$ inside
$\operatorname{Pic}^0_{Y/L}(L)$. Injectivity of the passage from actual
line bundles follows from Hilbert~90.

For completeness, let $x\in\bigcap_n nP$. Its image in $N$ vanishes,
since $N$ is finitely generated. Let $e$ annihilate the torsion subgroup
of $N$. For every $m\geq1$, choose $y\in P$ with $em y=x$.
Then the image of $y$ in $N$ is torsion, so $z=ey$ lies in $P^0$ and
$mz=x$. Thus $x$ is infinitely divisible in
$\operatorname{Pic}^0_{Y/L}(L)$, and (F47.H1) gives $x=0$. Finally, each
coordinate of a compatible torsion system in $T P$ is divisible by
every positive integer, so the first equality in (F47.H2) implies the
second. This proves the lemma, including the torsion in $N$.

\emph{Curves with prescribed points and transverse directions.}
We will use the following elementary consequence of Bertini. In a
smooth geometrically integral projective variety over an infinite
characteristic-zero field, a finite set of rational points lies on a
smooth geometrically integral complete-intersection curve. At any
specified point one may require its tangent line to be transverse to
a given smooth divisor. The construction is over the ground field;
one first makes the chosen points rational by a finite extension when
necessary. To check the imposed conditions, take sufficiently high
powers of an ample bundle. Serre vanishing makes evaluation on the
finite first-order neighbourhoods surjective. The linear systems
vanishing at the prescribed points embed the complement of those
points, and a general member is smooth there and geometrically
irreducible. Smoothness and the required independent differentials at
the finitely many prescribed points are nonempty open conditions on
the first jets. Intersect these opens and repeat until dimension one.
This is the usual prescribed-point version of Bertini; the needed
linear-system and irreducibility statements are Stacks, Lemmas
33.47.1, 33.47.3, and 37.32.3 \cite{FA_StacksBertini}.

\emph{Integral boundary residues and the global unit.}
Let $X/k$ be a connected regular affine scheme of finite type with
$x_0\in X(k)$, and let
\[
       \phi:X_{\mathrm{et}}\longrightarrow
                        (\mathbb G_{m,k})_{\mathrm{et}}
\]
be pinned over $B_k$. Write $\phi(a)$ for its value at a geometric
point. Over a finite extension $L/k$, a point $a\in X(L)$ has
$\phi(a)\in L^*$: the relative morphism of closed-point Galois topoi
forces the residue field of the image to embed in $L$. Base change of
$\phi$ to $L$ is obtained by slicing by the finite etale constant-field
object. These observations justify every evaluation used below.

The scheme $X$ is smooth and geometrically integral. In fact, in
characteristic zero regularity implies smoothness, and the algebraic
closure of $k$ in $k(X)$ is a finite field extension contained in all
normal local rings of $X$; evaluation at $x_0$ embeds it in $k$.
Choose a smooth projective compactification $Y$ of $X$ whose boundary
is a divisor with normal crossings, by characteristic-zero embedded
resolution. Concretely, take a projective closure, resolve it while
preserving the smooth open $X$, and principalize its boundary ideal;
the blowups are projective and supported outside $X$. The strong
resolution and monomialization theorem in
\cite[``The result'' and ``Idea of proof'', pp.~1--3]{FA_EncinasHauser2002}
supplies these operations. Write the boundary prime components over
$k$ as $D_1,\ldots,D_s$.
Only smooth points of individual boundary components will be used.

Pulling back the Kummer torsors of the coordinate $t$ gives
\[
       h=(h_n)_n\in\varprojlim_n H^1_{\mathrm{et}}(X,\mu_n).
                                                               \tag{F47.H3}
\]
Let $K=k(Y)$. The map
$H^1(X,\mu_n)\to K^*/K^{*n}$ is injective: a generically trivial
finite etale torsor on a normal integral scheme has a section, by
normalization. For every prime divisor $D_i$, the valuation of the
generic Kummer representative of $h_n$ defines a residue
$r_{i,n}\in\mathbb Z/n$. Compatibility gives
$r_i\in\widehat{\mathbb Z}$.

We claim $r_i\in\mathbb Z$. After a finite extension $L/k$, choose
an $L$-rational point $q$ on a geometric component above $D_i$, away
from the other boundary components, and a smooth complete-intersection
curve $\overline C\subset Y_L$ passing through $q$ transversely to
that component and through $x_0$. Its open part
$C=\overline C\cap X_L$ is an affine smooth geometrically integral
curve with an $L$-point. The curve calculation makes $h|_C$ the Kummer
class of an actual unit; its residue at $q$ is therefore an integer.
That residue is exactly $r_i$.

One can check this equality without an inverse-limit purity argument.
At $q$, the regular local ring $R$ is factorial and the only boundary
has one prime equation $z$. Hence
$\operatorname{Pic}(R[1/z])=0$ and
$R[1/z]^*=R^*z^{\mathbb Z}$. For each $n$, the restricted Kummer
class has a representative $u z^j$, with $u\in R^*$, and residue
$j\bmod n$. On the transverse curve $z$ is a local parameter and
$u$ remains a unit, giving the same residue for every $n$. The base
change from $D_i$ to a geometric component has multiplicity one in
characteristic zero. This proves the claim.

Set $D=\sum_i r_iD_i$, now an integral Cartier divisor on $Y$.
Its class lies in $\bigcap_n n\operatorname{Pic}(Y)$. Indeed, choose
a rational function $a_n\in K^*$ representing the generic class of
$h_n$. Its valuations at divisors inside $X$ are multiples of $n$, by
the Kummer description over a discrete valuation ring; at $D_i$ its
valuation is $r_i\bmod n$. Consequently
\[
                    \operatorname{div}(a_n)=D+nE_n
                                                               \tag{F47.H4}
\]
for an integral divisor $E_n$ on the regular scheme $Y$.
Thus $[D]=-n[E_n]$. By (F47.H2), $[D]=0$, and there is a rational
function $g\in k(Y)^*$ with $\operatorname{div}(g)=D$. In
particular, $g|_X$ is an actual regular unit.

Subtracting $[g]$ from $h$ leaves a compatible class unramified on
$Y$. Explicitly, the divisor of $a_n/g$ is divisible by $n$, so its
associated line bundle with an $n$th-power trivialization defines a
$\mu_n$-torsor on $Y$ extending $h_n-[g]$. Its uniqueness follows
from injectivity into $K^*/K^{*n}$, so these extensions are compatible
as $n$ varies. The Kummer exact sequences on $Y$ and (F47.H2) give
\[
  \varprojlim_n H^1(Y,\mu_n)=\widehat{k^*}.
                                                               \tag{F47.H5}
\]
To see that no derived-limit term is being suppressed, the maps
$k^*/k^{*mn}\to k^*/k^{*n}$ are surjective; alternatively, since the
compatible Picard-torsion class is zero, each class is represented
uniquely in $k^*/k^{*n}$, and these representatives are automatically
compatible. Here $\Gamma(Y,\mathcal O_Y)^*=k^*$ by geometric
integrality and properness.

Evaluation at $x_0$ identifies the constant class in (F47.H5) with the
actual element $c=\phi(x_0)/g(x_0)\in k^*$. Therefore
\[
        f=cg|_X\in\Gamma(X,\mathcal O_X)^*,\qquad
        h=[f],\qquad f(x_0)=\phi(x_0).                 \tag{F47.H6}
\]

\emph{Pointwise equality in every dimension.}
Let $a\in X(\overline k)$. Choose a finite extension $L/k$ making
$a$ rational and a smooth geometrically integral curve
$\overline C\subset Y_L$ through both $x_0$ and $a$; put
$C=\overline C\cap X_L$. In dimension one take $C=X_L$ itself.
The curve theorem gives a unit $f_C$ whose values equal those of
$\phi|_C$. Since both $f_C$ and $f|_C$ have Kummer class $h|_C$,
the unit $u=f_C/(f|_C)$ has zero Kummer class for every $n$. The
injection
\[
  \Gamma(C,\mathcal O_C)^*/\Gamma(C,\mathcal O_C)^{*n}
       \lhook\joinrel\longrightarrow H^1(C,\mu_n)
\]
shows that $u$ has an $n$th root among regular units for every $n$.
Its divisor on $\overline C$ is therefore divisible by every positive
integer and is zero. Thus $u\in L^*$ is constant. Evaluation at
$x_0$ gives $u=1$, and consequently
\[
                         \phi(a)=f(a)
                  \quad\hbox{for all }a\in X(\overline k).
                                                               \tag{F47.H7}
\]
In dimension zero $X=\operatorname{Spec}k$ and the assertion is
immediate. Notice that this final argument does not infer equality
from a finite collection of Kummer quotients: it uses the integral
curve realization and equality in every Kummer degree. It also does
not require the residual roots of unity to form a chosen compatible
system.

\emph{Recovery of the entire geometric morphism.}
The assertion just proved holds for every connected regular affine
finite-type source with a rational point, and after every finite
extension of the ground field. It is exactly condition (4.18.5) of
Carlson--Haine--Wolf, Theorem~4.18 \cite{CHW2024}. Their equivalence
with (4.18.1) therefore gives the full pinned reconstruction statement
for this field class, with absolute weak normalization as in their
definition. This application addresses all etale objects, not just
the point map of one morphism.

Here is the categorical part of that implication. Two copies of
$\mathbb G_m$ cover $\mathbb A^1$, so the pointwise units just
constructed give pointwise algebraic maps to $\mathbb A^1$. Affine
coordinates and equality on geometric points of a reduced source
then reconstruct the addition and multiplication relations and hence
a map to every affine target. Apply this construction also to each
sliced pinned morphism over a quasicompact separated etale target
object $V$. Its inverse image is again represented by a
quasicompact separated etale source scheme, by CHW Corollary~A.9;
one works on its affine components and, when needed, after a finite
extension supplying a rational point. Uniqueness on geometric points
glues these maps and makes them compatible with all morphisms of
etale objects. Proposition~4.14 then identifies the original inverse
image functor with the functor of the reconstructed scheme map.
Finally, the $h$-descent part of Theorem~4.18 handles arbitrary
finite-type sources. Thus no openness assertion for an unconstructed
birational Galois homomorphism, and no inference from finite set
coefficients alone, is needed.

\end{proof}
